% GENERATED FILE. Edit canonical lessons, then run npm run generate:books.
% canonical-source-hash: fnv1a64:4ac6573f798fec5d
% canonical-lessons: UR-C268-lagna, UR-C268-khayal, UR-C268-ghalat, UR-W268-ghain, UR-C268-mujhe-lagta-hai, UR-C268-ki-vajah-se, UR-R268-opinion-and-reason

\chapter{What I Think, and Why}
\label{ch:ur-what-i-think-and-why}
\hlchaptermodality{268}

\begin{chapteropening}
\textbf{By the end of this chapter:} \emph{I can say what I think (mujhe lagtā hai ki …, mere khayāl meṅ …) and give a reason (kyūṅki …, … kī vajah se).}

Complete UR-R268-opinion-and-reason: say what you think, and why.
\end{chapteropening}

\section[\texorpdfstring{\ur{لگنا} (lagnā)}{لگنا (lagnā)}]{\ur{لگنا} (lagnā) — to seem, to be felt}

\label{lesson:UR-C268-lagna}



\begin{quote}
Before the new one: say \textquotedblleft{}hunger\textquotedblright{} and \textquotedblleft{}thirst\textquotedblright{}. (\emph{bhūk}, \emph{pyās}.) Urdu does not say \textquotedblleft{}I am hungry\textquotedblright{}. It says that hunger has struck you.
\end{quote}

\subsection*{You'll want to know: \ur{لگنا}}
\textbf{\ur{لگنا}} — \emph{lagnā} — \textquotedblleft{}to seem; to be felt\textquotedblright{}.

\begin{quote}
\textbf{\ur{مجھے} \ur{بھوک} \ur{لگی} \ur{ہے۔}}
\end{quote}

\begin{quote}
\emph{mujhe bhūk lagī hai.} — \textquotedblleft{}I'm hungry.\textquotedblright{}
\end{quote}

The one who feels it takes \textbf{\ur{مجھے}} (\emph{mujhe}), and the verb follows the thing felt: \textbf{\ur{بھوک}} (\emph{bhūk}) is feminine, so \textbf{\ur{لگی}} (\emph{lagī}).

\begin{quote}
\textbf{\ur{مجھے} \ur{یہ} \ur{فلم} \ur{اچھی} \ur{لگی۔}}
\end{quote}

\begin{quote}
\emph{mujhe yih film achchhī lagī.} — \textquotedblleft{}I liked this film.\textquotedblright{}
\end{quote}

Word for word, \textquotedblleft{}this film seemed good to me\textquotedblright{}. \textbf{\ur{فلم}} is feminine: \textbf{\ur{اچھی} \ur{لگی}}.

\begin{cousinweb}
\textbf{\ur{لگنا}} (\emph{lagnā}) is inherited. CDIAL traces it to Sanskrit \emph{lagyati}, \textquotedblleft{}adheres to\textquotedblright{} (CDIAL 10895): a hunger or a liking attaches itself to the one who feels it.
\end{cousinweb}

\subsection*{Your turn}
Say these aloud:

\begin{itemize}
\raggedright
  \item \emph{lagnā}, then \emph{mujhe bhūk lagī hai}
  \item \textquotedblleft{}I'm thirsty\textquotedblright{} — \emph{mujhe pyās lagī hai}
  \item \textquotedblleft{}I liked this film\textquotedblright{} — \emph{mujhe yih film achchhī lagī}
\end{itemize}

\begin{tcolorbox}[breakable,colback=teal!4,colframe=teal!35!black,title={Before you move on}]
What does \textbf{\ur{لگنا}} mean? (\textquotedblleft{}To seem, to be felt\textquotedblright{}.) \textquotedblleft{}I'm hungry\textquotedblright{}? (\textbf{\ur{مجھے} \ur{بھوک} \ur{لگی} \ur{ہے}}.) Why \textbf{\ur{لگی}}? (\textbf{\ur{بھوک}} is feminine.)
\end{tcolorbox}
\section[\texorpdfstring{\ur{خیال} (khayāl)}{خیال (khayāl)}]{\ur{خیال} (khayāl) — a thought, an idea}

\label{lesson:UR-C268-khayal}



\begin{quote}
Before the new one: say \textquotedblleft{}I liked it\textquotedblright{}, of a film. (\emph{mujhe achchhī lagī}.) Then \textquotedblleft{}to think\textquotedblright{}. (\emph{sochnā}.)
\end{quote}

\subsection*{You'll want to know: \ur{خیال}}
\textbf{\ur{خیال}} — \emph{khayāl} — \textquotedblleft{}a thought; an idea\textquotedblright{}. It is masculine.

\begin{quote}
\textbf{\ur{یہ} \ur{اچھا} \ur{خیال} \ur{ہے۔}}
\end{quote}

\begin{quote}
\emph{yih achchhā khayāl hai.} — \textquotedblleft{}That's a good idea.\textquotedblright{}
\end{quote}

\begin{quote}
\textbf{\ur{مجھے} \ur{ایک} \ur{خیال} \ur{آیا۔}}
\end{quote}

\begin{quote}
\emph{mujhe ek khayāl āyā.} — \textquotedblleft{}An idea came to me.\textquotedblright{}
\end{quote}

An idea comes to you, as hunger strikes you: \textbf{\ur{مجھے}}, and \textbf{\ur{آیا}} (\emph{āyā}) for a masculine thought.

\begin{cousinweb}
\textbf{\ur{خیال}} (\emph{khayāl}) is usually traced, through Persian, to Arabic \emph{khayāl}, \textquotedblleft{}a fancy, an image in the mind\textquotedblright{}.
\end{cousinweb}

\subsection*{Your turn}
Say these aloud:

\begin{itemize}
\raggedright
  \item \emph{khayāl}, then \emph{yih achchhā khayāl hai}
  \item \textquotedblleft{}That's a bad idea\textquotedblright{} — \emph{yih burā khayāl hai}
  \item \textquotedblleft{}An idea came to me\textquotedblright{} — \emph{mujhe ek khayāl āyā}
\end{itemize}

\begin{tcolorbox}[breakable,colback=teal!4,colframe=teal!35!black,title={Before you move on}]
What does \textbf{\ur{خیال}} mean? (\textquotedblleft{}A thought, an idea\textquotedblright{}.) Masculine or feminine? (Masculine: \textbf{\ur{اچھا} \ur{خیال}}.)
\end{tcolorbox}
\section[\texorpdfstring{\ur{غلط} (ghalat)}{غلط (ghalat)}]{\ur{غلط} (ghalat) — wrong}

\label{lesson:UR-C268-ghalat}



\begin{quote}
Before the new one: say \textquotedblleft{}a good idea\textquotedblright{}. (\emph{achchhā khayāl}.) Then \textquotedblleft{}fine, correct\textquotedblright{}. (\emph{ṭhīk}.) This lesson gives its opposite.
\end{quote}

\subsection*{You'll want to know: \ur{غلط}}
\textbf{\ur{غلط}} — \emph{ghalat} — \textquotedblleft{}wrong\textquotedblright{}.

\begin{quote}
\textbf{\ur{یہ} \ur{جواب} \ur{غلط} \ur{ہے۔}}
\end{quote}

\begin{quote}
\emph{yih javāb ghalat hai.} — \textquotedblleft{}This answer is wrong.\textquotedblright{}
\end{quote}

\begin{quote}
\textbf{\ur{میں} \ur{غلط} \ur{تھا۔}}
\end{quote}

\begin{quote}
\emph{maiṅ ghalat thā.} — \textquotedblleft{}I was wrong.\textquotedblright{}
\end{quote}

It does not change for a man or a woman: \textbf{\ur{میں} \ur{غلط} \ur{تھی}} (\emph{maiṅ ghalat thī}).

\begin{cousinweb}
\textbf{\ur{غلط}} (\emph{ghalat}) is usually traced, through Persian, to Arabic \emph{ghalaṭ}, \textquotedblleft{}an error\textquotedblright{}. Its first letter is new to your hand; the next lesson takes it out.
\end{cousinweb}

\subsection*{Your turn}
Say these aloud:

\begin{itemize}
\raggedright
  \item \emph{ghalat}, then \emph{yih javāb ghalat hai}
  \item \textquotedblleft{}This answer is right\textquotedblright{} — \emph{yih javāb ṭhīk hai}
  \item \textquotedblleft{}I was wrong\textquotedblright{}, as a woman — \emph{maiṅ ghalat thī}
\end{itemize}

\begin{tcolorbox}[breakable,colback=teal!4,colframe=teal!35!black,title={Before you move on}]
What does \textbf{\ur{غلط}} mean? (\textquotedblleft{}Wrong\textquotedblright{}.) Its opposite? (\textbf{\ur{ٹھیک}}.)
\end{tcolorbox}
\section[\texorpdfstring{\ur{غ} (ghain)}{غ (ghain)}]{\ur{غ} (ghain) — the letter that opens \ur{غلط}}

\label{lesson:UR-W268-ghain}



\begin{quote}
Before the new one: write \ur{ع} once, and name it. (\emph{ain}.)

Now say \textbf{\ur{غلط}} (\emph{ghalat}), \textquotedblleft{}wrong\textquotedblright{}. This lesson takes one letter out of it and puts it in your hand.
\end{quote}

\subsection*{Script you'll notice: \ur{غ}}
\textbf{\ur{غ}} — \emph{ghain}, the sound \emph{gh}, made far back in the throat.

It opens \textbf{\ur{غلط}} (\emph{ghalat}). You already write \textbf{\ur{ع}}: \textbf{\ur{غ}} is that shape with \textbf{one dot} above it. At the start of a word, as here, it is the small open head alone, joined on to the letter after it; standing alone it keeps the deep bowl below the line.

\subsection*{Writing: \ur{غ}}
Put your pen on \ur{غ} and follow its line. Copy the shape you can see — slowly, and larger than it is printed.

\begin{quote}
This book does not yet tell you \textbf{where to start this letter or which way to travel}. It is not written down here until it can be written down with a source. Copying what is in front of you needs no such source, and it is how the shape gets into your hand in the meantime.
\end{quote}

\subsection*{Your turn}
\begin{itemize}
\raggedright
  \item \emph{Look:} at \textbf{\ur{غلط}} and point at \ur{غ} inside it
  \item \emph{Trace it:} \ur{غ} three times, saying \emph{gh} as you finish each one
  \item \emph{Write it:} \textbf{\ur{غلط}} once, joined, with the dot above its first letter
  \item \emph{Read it:} \textbf{\ur{غلط}} aloud once more
\end{itemize}

\begin{tcolorbox}[breakable,colback=teal!4,colframe=teal!35!black,title={Before you move on}]
Which letter is this — \ur{غ}? (\textbf{ghain}, \emph{gh}.) What tells it from \ur{ع}? (The dot.) Which word did it come from? (\textbf{\ur{غلط}}, \emph{ghalat}, \textquotedblleft{}wrong\textquotedblright{}.)
\end{tcolorbox}
\section[\texorpdfstring{\ur{مجھے} \ur{لگتا} \ur{ہے} (mujhe lagtā hai)}{مجھے لگتا ہے (mujhe lagtā hai)}]{\ur{مجھے} \ur{لگتا} \ur{ہے} (mujhe lagtā hai) — what I think}

\label{lesson:UR-C268-mujhe-lagta-hai}



\begin{quote}
Say \textquotedblleft{}I'm hungry\textquotedblright{}. (\emph{mujhe bhūk lagī hai}.) Then \textquotedblleft{}it will rain\textquotedblright{}. (\emph{bārish ho gī}.) Join the two verbs, and you have an opinion.
\end{quote}

\subsection*{You'll want to know: \ur{مجھے} \ur{لگتا} \ur{ہے} \ur{کہ}}
\begin{quote}
\textbf{\ur{مجھے} \ur{لگتا} \ur{ہے۔}}
\end{quote}

\begin{quote}
\emph{mujhe lagtā hai.} — \textquotedblleft{}I think so.\textquotedblright{}
\end{quote}

Word for word, \textquotedblleft{}it seems to me\textquotedblright{}. \textbf{\ur{کہ}} (\emph{ki}), \textquotedblleft{}that\textquotedblright{}, hangs the thought on it:

\begin{quote}
\textbf{\ur{مجھے} \ur{لگتا} \ur{ہے} \ur{کہ} \ur{بارش} \ur{ہو} \ur{گی۔}}
\end{quote}

\begin{quote}
\emph{mujhe lagtā hai ki bārish ho gī.} — \textquotedblleft{}I think it will rain.\textquotedblright{}
\end{quote}

\begin{quote}
\textbf{\ur{مجھے} \ur{لگتا} \ur{ہے} \ur{کہ} \ur{وہ} \ur{نہیں} \ur{آئے} \ur{گا۔}}
\end{quote}

\begin{quote}
\emph{mujhe lagtā hai ki voh nahīṅ āe gā.} — \textquotedblleft{}I don't think he'll come.\textquotedblright{}
\end{quote}

\subsection*{You'll want to know: in my view}
\begin{quote}
\textbf{\ur{میرے} \ur{خیال} \ur{میں} \ur{یہ} \ur{ٹھیک} \ur{ہے۔}}
\end{quote}

\begin{quote}
\emph{mere khayāl meṅ yih ṭhīk hai.} — \textquotedblleft{}In my view, this is right.\textquotedblright{}
\end{quote}

\begin{quote}
\textbf{\ur{میری} \ur{رائے} \ur{میں} \ur{یہ} \ur{غلط} \ur{ہے۔}}
\end{quote}

\begin{quote}
\emph{merī rāe meṅ yih ghalat hai.} — \textquotedblleft{}In my opinion, this is wrong.\textquotedblright{}
\end{quote}

\textbf{\ur{خیال}} is masculine, so \textbf{\ur{میرے}} (\emph{mere}); \textbf{\ur{رائے}} (\emph{rāe}) is feminine, so \textbf{\ur{میری}} (\emph{merī}). To ask:

\begin{quote}
\textbf{\ur{آپ} \ur{کا} \ur{کیا} \ur{خیال} \ur{ہے؟}}
\end{quote}

\begin{quote}
\emph{āp kā kyā khayāl hai?} — \textquotedblleft{}What do you think?\textquotedblright{}
\end{quote}

\subsection*{Your turn}
Say these aloud:

\begin{itemize}
\raggedright
  \item \textquotedblleft{}I think it will rain\textquotedblright{} — \emph{mujhe lagtā hai ki bārish ho gī}
  \item \textquotedblleft{}I think she will come\textquotedblright{} — \emph{mujhe lagtā hai ki voh āe gī}
  \item \textquotedblleft{}I think they will come\textquotedblright{} — \emph{mujhe lagtā hai ki voh āeṅ ge}
\end{itemize}

\begin{tcolorbox}[breakable,colback=teal!4,colframe=teal!35!black,title={Before you move on}]
\textquotedblleft{}I think so\textquotedblright{}? (\textbf{\ur{مجھے} \ur{لگتا} \ur{ہے}}.) \textquotedblleft{}In my view\textquotedblright{}? (\textbf{\ur{میرے} \ur{خیال} \ur{میں}}.) Why \textbf{\ur{میری} \ur{رائے}}? (\textbf{\ur{رائے}} is feminine.)
\end{tcolorbox}
\section[\texorpdfstring{\ur{کی} \ur{وجہ} \ur{سے} (kī vajah se)}{کی وجہ سے (kī vajah se)}]{\ur{کی} \ur{وجہ} \ur{سے} (kī vajah se) — the reason}

\label{lesson:UR-C268-ki-vajah-se}



\begin{quote}
Say \textquotedblleft{}because\textquotedblright{} and \textquotedblleft{}a reason\textquotedblright{}. (\emph{kyūṅki}, \emph{vajah}.) Then \textquotedblleft{}I think so\textquotedblright{}. (\emph{mujhe lagtā hai}.) An opinion wants a reason.
\end{quote}

\subsection*{You'll want to know: \ur{کیونکہ} and \ur{چونکہ}}
\begin{quote}
\textbf{\ur{میں} \ur{نہیں} \ur{آیا۔}}
\end{quote}

\begin{quote}
\emph{maiṅ nahīṅ āyā.} — \textquotedblleft{}I didn't come.\textquotedblright{}
\end{quote}

\begin{quote}
\textbf{\ur{میں} \ur{نہیں} \ur{آیا} \ur{کیونکہ} \ur{میں} \ur{بیمار} \ur{تھا۔}}
\end{quote}

\begin{quote}
\emph{maiṅ nahīṅ āyā kyūṅki maiṅ bīmār thā.} — \textquotedblleft{}I didn't come because I was ill.\textquotedblright{}
\end{quote}

\textbf{\ur{کیونکہ}} (\emph{kyūṅki}) puts the reason second. \textbf{\ur{چونکہ}} (\emph{chūṅki}) puts it first, and the other half opens with \textbf{\ur{اس} \ur{لیے}} (\emph{is liye}):

\begin{quote}
\textbf{\ur{چونکہ} \ur{میں} \ur{بیمار} \ur{تھا،}}
\end{quote}

\begin{quote}
\emph{chūṅki maiṅ bīmār thā,}
\end{quote}

\begin{quote}
\textbf{\ur{اس} \ur{لیے} \ur{میں} \ur{نہیں} \ur{آیا۔}}
\end{quote}

\begin{quote}
\emph{is liye maiṅ nahīṅ āyā.} — \textquotedblleft{}Since I was ill, I didn't come.\textquotedblright{}
\end{quote}

\subsection*{You'll want to know: a noun as the reason}
\begin{quote}
\textbf{\ur{بارش} \ur{کی} \ur{وجہ} \ur{سے} \ur{ہم} \ur{گھر} \ur{پر} \ur{رہے۔}}
\end{quote}

\begin{quote}
\emph{bārish kī vajah se ham ghar par rahe.} — \textquotedblleft{}Because of the rain we stayed at home.\textquotedblright{}
\end{quote}

\textbf{\ur{وجہ}} (\emph{vajah}) is feminine, so \textbf{\ur{کی}} (\emph{kī}). The noun goes in front:

\begin{quote}
\textbf{\ur{دیر} \ur{کی} \ur{وجہ} \ur{سے} \ur{وہ} \ur{ناراض} \ur{تھی۔}}
\end{quote}

\begin{quote}
\emph{der kī vajah se voh nārāz thī.} — \textquotedblleft{}She was cross because of the delay.\textquotedblright{}
\end{quote}

\subsection*{Your turn}
Say these aloud:

\begin{itemize}
\raggedright
  \item \textquotedblleft{}Because of the rain I stayed at home\textquotedblright{} — \emph{bārish kī vajah se maiṅ ghar par rahā}
  \item \textquotedblleft{}Because of the delay I stayed at home\textquotedblright{} — \emph{der kī vajah se maiṅ ghar par rahā}
  \item \textquotedblleft{}Because of the heat I stayed at home\textquotedblright{} — \emph{garmī kī vajah se maiṅ ghar par rahā}
\end{itemize}

\begin{tcolorbox}[breakable,colback=teal!4,colframe=teal!35!black,title={Before you move on}]
\textquotedblleft{}Because I was ill\textquotedblright{}? (\textbf{\ur{کیونکہ} \ur{میں} \ur{بیمار} \ur{تھا}}.) \textquotedblleft{}Because of the rain\textquotedblright{}? (\textbf{\ur{بارش} \ur{کی} \ur{وجہ} \ur{سے}}.) Why \textbf{\ur{کی}}? (\textbf{\ur{وجہ}} is feminine.)
\end{tcolorbox}
\section[(review)]{Review — what I think, and why}

\label{lesson:UR-R268-opinion-and-reason}



\begin{quote}
Say \textquotedblleft{}having eaten\textquotedblright{} and \textquotedblleft{}after the meal\textquotedblright{}. (\emph{khā kar}, \emph{khāne ke bād}.)
\end{quote}

\subsection*{Your turn}
\begin{itemize}
\raggedright
  \item \emph{Say it:} I'm hungry. I liked this film.
  \item \emph{Say it:} I think it will rain. What do you think?
  \item \emph{Say it:} In my view this is right. In my opinion that is wrong.
  \item \emph{Say it:} I didn't come because I was ill.
  \item \emph{Say it:} Because of the rain we stayed at home.
  \item \emph{Write it:} \textbf{\ur{غلط}} once, and point at \ur{غ} inside it
\end{itemize}

\begin{tcolorbox}[breakable,colback=teal!4,colframe=teal!35!black,title={Before you move on}]
\textquotedblleft{}I think he won't come\textquotedblright{}? (\textbf{\ur{مجھے} \ur{لگتا} \ur{ہے} \ur{کہ} \ur{وہ} \ur{نہیں} \ur{آئے} \ur{گا}}.) Now the other way round: \textbf{\ur{چونکہ}} or \textbf{\ur{کیونکہ}} to put the reason first? (\textbf{\ur{چونکہ}}.) \textquotedblleft{}Because of the delay\textquotedblright{}? (\textbf{\ur{دیر} \ur{کی} \ur{وجہ} \ur{سے}}.)
\end{tcolorbox}
